\documentclass[11pt]{article}

\usepackage[margin=1in]{geometry}
\usepackage{amsmath,amssymb,amsfonts,mathtools}
\usepackage{bm}
\usepackage{graphicx}
\usepackage[font=small]{caption}
\usepackage{subcaption}
\usepackage{titlesec}
\usepackage[numbers,comma,sort&compress]{natbib}
\usepackage{authblk}
\usepackage[colorlinks]{hyperref}
\hypersetup{citecolor=blue,linkcolor=blue,urlcolor=blue}

\renewcommand{\thefigure}{S\arabic{figure}}
\renewcommand{\theequation}{S\arabic{equation}}
\renewcommand{\thesection}{Supplementary Note \arabic{section}}
\titleformat{\section}{\large\bfseries}{\thesection}{0.75em}{}

\newcommand{\mc}{\mathcal}
\newcommand{\eps}{\varepsilon}
\newcommand{\Lind}{\mathcal{L}}
\newcommand{\Tr}{\operatorname{Tr}}
\newcommand{\ad}{\operatorname{ad}}
\newcommand{\comm}{\operatorname{comm}}
\newcommand{\supp}{\operatorname{supp}}

\title{Supplementary Information for\\
Joint Mitigation of Algorithmic and Physical Errors in Noisy Hamiltonian Simulation}
\author{}
\date{}

\begin{document}

\maketitle

This Supplementary Information collects the technical and experimental material that is deliberately kept out of the short Nature Physics main text. Supplementary Notes 1 and 2 give the extrapolation and finite-order BCH details behind the coupled path. Supplementary Notes 3--6 document the hardware circuit, readout calibration, learned sparse Pauli--Lindblad channel, hardware extrapolation data and numerical validation.

The hardware data are processed through the following pipeline. Raw computational-basis bitstrings are first converted into readout-mitigated estimates of the average magnetization. Cliffordized training circuits are then used to learn a layerwise sparse Pauli--Lindblad model for the compiled device circuit. The learned model determines random Pauli insertions that realize several amplified noise levels. For each fixed Trotter step number, these data are extrapolated to zero physical noise. Finally, the zero-noise estimates at different Trotter step numbers are extrapolated in $x=1/N_{\mathrm{Trotter}}$ to reduce the deterministic product-formula error.

\section{Coupled Richardson Path}

The main text considers expectation values of the form
\begin{equation}
    f(s)=g(s,\lambda(s)),\qquad
    \lambda(s)=c(sT)^{p+1},
\end{equation}
where $s$ is the product-formula step size, $T$ is the simulated time, $p$ is the product-formula order and $\lambda$ is the effective noise strength. The target expectation value is $f(0)$.

Suppose that, over the range of step sizes used in the extrapolation,
\begin{equation}
    f(s)=\sum_{\ell=0}^{m-1} a_\ell s^\ell + R_m(s),
\end{equation}
where $R_m$ denotes the remainder after the first $m$ powers. Given nodes $s_j=T/r_j$, the Richardson estimate is
\begin{equation}
    F_m=\sum_{j=1}^{m} b_j f(s_j),\qquad
    \sum_{j=1}^{m} b_j s_j^\ell =
    \begin{cases}
    1, & \ell=0,\\
    0, & 1\le \ell\le m-1 .
    \end{cases}
\end{equation}
Thus
\begin{equation}
    |F_m-f(0)|\le \|\bm b\|_1\max_j |R_m(s_j)|.
\end{equation}
The role of the coupled path is to make the leading physical-noise bias and the leading Trotter bias appear in the same polynomial expansion in $s$. For a $p$-th order product formula, the deterministic Trotter term begins at order $s^{p+1}$. Since the leading Markovian noise contribution is linear in $\lambda$, the choice $\lambda(s)=c(sT)^{p+1}$ places both contributions at the same order. A smaller exponent would make noise dominate, while a larger exponent would leave the Trotter term as the bottleneck.

For comparison, a tensor-product two-dimensional extrapolation with coefficients $b_jb'_\ell$ has coefficient one-norm $\|\bm b\|_1\|\bm b'\|_1$. Thus every grid point must be estimated more accurately by the extra factor $\|\bm b'\|_1$ compared with a one-dimensional path using the same step-size nodes. For generic noise-direction nodes $x_\ell$, the coefficient norm is bounded by the Lagrange expression
\begin{equation}
    \|\bm b'\|_1
    =
    \sum_{\ell}
    \left|
    \prod_{q\ne \ell}\frac{x_q}{x_q-x_\ell}
    \right|.
\end{equation}
If this direction is not well conditioned, the resulting sampling overhead can dominate the apparent simplicity of a rectangular grid.

\section{Finite-Order BCH Remainder}

For one noisy product-formula step, define
\begin{equation}
    \mathcal{K}_{2v-1}=-i sT a_v \ad_{H_{\gamma_v}},
    \qquad
    \mathcal{K}_{2v}=c(sT)^{p+1}\Lind_v .
\end{equation}
The single-step channel is a product of exponentials of these generators. Formally,
\begin{equation}
    \prod_v e^{\mathcal{K}_v}
    =
    \exp\left(\sum_{q=1}^{\infty}\Phi_q\right),
\end{equation}
where $\Phi_q$ is a homogeneous sum of $q$-fold nested commutators. For local Hamiltonians, the full BCH series need not converge for any nonzero step size because high-order nested commutator bounds may grow factorially. We therefore use only the truncated generator
\begin{equation}
    \mathcal{G}_{(q_0)}(s)=\frac{1}{sT}\sum_{q=1}^{q_0}\Phi_q(sT,c(sT)^{p+1}).
\end{equation}

The truncation error is controlled by doubly right-nested commutator sums
\begin{equation}
    \bar{\alpha}_{\comm}^{(q_1,\ldots,q_d)}
    =
    \sum
    \left\|
    \left[
      [\mathcal{K}_{v_{1,1}},\ldots,\mathcal{K}_{v_{1,q_1}}],
      \ldots,
      [\mathcal{K}_{v_{d,1}},\ldots,\mathcal{K}_{v_{d,q_d}}]
    \right]
    \right\|.
\end{equation}
For $k$-local, $g$-extensive superoperators, the support of a nested commutator grows only when a new local term overlaps the existing support. As a result, these sums scale with local factors such as $kg$ and support-growth factors rather than with the global norm of the full Hamiltonian. Choosing $q_0=O(\log(N/\eps))$ makes the finite BCH remainder smaller than the target accuracy while preserving the commutator-sensitive scaling of product formulas.

\section{Circuit Mapping and Readout Calibration}

The hardware experiment uses a connected 9-qubit subgraph of a 13-qubit superconducting processor. The simulated observable is the average longitudinal magnetization,
\begin{equation}
    O=\frac{1}{n}\sum_{i=1}^{n} Z_i .
\end{equation}
Each second-order symmetric Trotter step consists of transverse-field rotations and an Ising interaction layer. The interaction graph is decomposed into disjoint $R_{ZZ}$ patterns so that gates within one pattern are executed in parallel. The main text shows the circuit and topology; this note records the calibration data used before the learned-noise step.

Before learning the effective gate-noise channel, we mitigate readout error at the observable level. Supplementary Fig.~\ref{fig:supp-readout} shows both the full confusion matrix on the 9-qubit subgraph and the local single-qubit calibration matrices. The global matrix is strongly diagonally dominant, indicating that the readout channel is close to an independent bit-flip model on the relevant device subgraph. The local matrices capture the dominant contribution for the magnetization observable while avoiding the sampling cost of inverting the full global matrix in every data-processing step.

\begin{figure}[!htbp]
    \centering
    \begin{subfigure}{0.48\linewidth}
        \centering
        \includegraphics[width=\linewidth]{figures/supple/supp_readout_global.pdf}
        \caption{}
    \end{subfigure}
    \begin{subfigure}{0.48\linewidth}
        \centering
        \includegraphics[width=\linewidth]{figures/supple/supp_readout_local.pdf}
        \caption{}
    \end{subfigure}
    \caption{
    Readout calibration data for the selected 9-qubit device subgraph.
    (a) Full $2^9\times2^9$ global confusion matrix, whose entries give the conditional probabilities of measured bitstrings given prepared computational-basis states.
    (b) Single-qubit local calibration matrices used for observable-level readout mitigation of the magnetization.
    The diagonal dominance of the calibration data supports local observable-level correction for the measured average magnetization.
    }
    \label{fig:supp-readout}
\end{figure}

\section{Learned Sparse Pauli--Lindblad Channel}

The learned-channel quantum error mitigation workflow assumes an effective sparse Pauli--Lindblad model for each relevant circuit layer~\cite{berg_probabilistic_2023},
\begin{equation}
    \Lind(\rho)=\sum_{k\in\mathcal{K}}\lambda_k(P_k\rho P_k-\rho),
\end{equation}
where $\mathcal{K}$ contains the Pauli components retained by the sparse model. The corresponding channel can be written as a product of Pauli dephasing components,
\begin{equation}
    \Lambda(\rho)=
    \prod_{k\in\mathcal{K}}
    \left(w_k \rho +(1-w_k)P_k\rho P_k\right),
    \qquad
    w_k=\frac{1+e^{-2\lambda_k}}{2}.
\end{equation}

Training circuits are produced by replacing the target circuit rotation angles by Clifford angles in $\{0,\pi/2,\pi,3\pi/2\}$ while preserving the same circuit layout, qubit mapping, layer grouping and native-gate compilation. Their noiseless expectation values are computed classically, while their noisy expectation values are measured on hardware after readout mitigation. Solving the resulting sparse inverse problem yields the effective Pauli--Lindblad rates. The learned channel is an effective model for the observable and circuit family: components that are invisible to the measured expectation or degenerate with other components need not be separately identified. This is why the main text treats the model as a calibrated noise knob rather than as microscopic noise tomography.

Supplementary Fig.~\ref{fig:supp-learned-hardware-noise} shows a representative learned hardware noise channel for a Trotterized circuit with step number $N=3$. The figure displays the inferred Pauli-noise components across the Clifford layers of the compiled circuit, giving a layer-resolved view of the channel used to generate controlled noise-gain factors for ZNE.

\begin{figure}[!htbp]
    \centering
    \includegraphics[width=0.9\linewidth,height=0.84\textheight,keepaspectratio]{figures/supple/supp_learned_hardware_noise_model.pdf}
    \caption{
    Learned effective hardware noise model for a Trotterized Ising circuit with Trotter step number $N=3$.
    The plot shows the inferred sparse Pauli--Lindblad components associated with the Clifford layers in the compiled hardware circuit.
    These learned components provide the calibrated noise channel used for controlled noise amplification in the subsequent ZNE step.
    }
    \label{fig:supp-learned-hardware-noise}
\end{figure}

\section{Hardware ZNE and Trotter Extrapolation}

After learning the effective noise channel, we use random Pauli insertions to implement gain factors for the noise strength. For each fixed Trotter step number, the amplified-noise expectation values are fitted as functions of the gain factor and extrapolated to the zero-noise limit. In the hardware data we compare low-order polynomial fits with exponential fitting families; the latter are natural for Pauli--Lindblad channels because the channel eigenvalues depend exponentially on the learned rates. Supplementary Fig.~\ref{fig:supp-zne} shows hardware ZNE data for two representative evolution settings. In both cases, the measured expectation values vary smoothly with the gain factor, and the zero-noise estimates are consistent with the corresponding noiseless simulation references within the accuracy of the fit.

\begin{figure}[!htbp]
    \centering
    \begin{subfigure}{0.88\linewidth}
        \centering
        \includegraphics[width=\linewidth]{figures/archive/supp_hardware_zne_minus1.pdf}
        \caption{}
    \end{subfigure}
    \vspace{0.8em}
    \begin{subfigure}{0.88\linewidth}
        \centering
        \includegraphics[width=\linewidth]{figures/archive/supp_hardware_zne_minus2.pdf}
        \caption{}
    \end{subfigure}
    \caption{
    Hardware zero-noise extrapolation using the learned Pauli--Lindblad channel.
    (a) Evolution setting $r_xT=r_{zz}T=-1$.
    (b) Evolution setting $r_xT=r_{zz}T=-2$.
    The extrapolated value at zero noise is obtained by fitting the measured dependence on the gain factor.
    These data provide the finite-Trotter, zero-noise estimates that are subsequently used for Trotter-step extrapolation.
    }
    \label{fig:supp-zne}
\end{figure}

The hardware ZNE stage removes the physical-noise bias at fixed Trotter step number. It does not remove the deterministic error from the finite-step product formula. For this reason, the ZNE-corrected estimates are subsequently extrapolated in $x=1/N_{\mathrm{Trotter}}$. For the second-order symmetric formula, the leading product-formula corrections have an even-power structure in $x$. Supplementary Fig.~\ref{fig:supp-trotter} shows the corresponding even-polynomial Trotter extrapolation and the residual error after extrapolation. The zero-step estimates are close to the independently computed non-Trotter reference.

\begin{figure}[!htbp]
    \centering
    \begin{subfigure}{0.48\linewidth}
        \centering
        \includegraphics[width=\linewidth]{figures/archive/supp_trotter_extrapolation_minus2.pdf}
        \caption{}
    \end{subfigure}
    \begin{subfigure}{0.48\linewidth}
        \centering
        \includegraphics[width=\linewidth]{figures/archive/supp_trotter_error_residual_minus2.pdf}
        \caption{}
    \end{subfigure}
    \caption{
    Trotter-step extrapolation after noise mitigation for the setting $r_xT=r_{zz}T=-2$.
    (a) Extrapolation of the corrected expectation value as a function of $x=1/N_{\mathrm{Trotter}}$ using even-polynomial fits.
    (b) Comparison between finite-step Trotter errors and the residual errors after zero-step extrapolation.
    }
    \label{fig:supp-trotter}
\end{figure}

\section{Numerical Learned-Channel Validation}

To check that the learning procedure scales beyond the 9-qubit hardware instance, we also use classically simulated noisy circuits. The simulations use the same Trotterized Ising structure and average magnetization observable as in the hardware experiments. A synthetic sparse Pauli noise model is applied to the circuit, and LCQEM is run using finite-shot Clifford training data. The reference tensor-network simulations in the main text use \texttt{quimb}~\cite{Gray2018}.

Supplementary Fig.~\ref{fig:supp-numerical-learning} shows a representative learned-noise validation on a $5\times6$ lattice. The inferred noise parameters reproduce the noisy target-circuit expectation values under several gain factors. This comparison isolates the learning step: the ``realistic'' noisy expectations are generated directly from the synthetic noise model, whereas the learned expectations are generated from the fitted sparse Pauli--Lindblad channel.

\begin{figure}[!htbp]
    \centering
    \begin{subfigure}{0.48\linewidth}
        \centering
        \includegraphics[width=\linewidth]{figures/supple/supp_numerical_learned_noise_params.pdf}
        \caption{}
    \end{subfigure}
    \begin{subfigure}{0.48\linewidth}
        \centering
        \includegraphics[width=\linewidth]{figures/supple/supp_numerical_learned_expectations_5x6.pdf}
        \caption{}
    \end{subfigure}
    \caption{
    Numerical validation of the learned sparse Pauli--Lindblad channel.
    (a) Learned noise-channel parameters for a representative $5\times6$ Trotter circuit instance.
    (b) Comparison between noisy expectation values obtained directly from the synthetic noise model and those generated from the learned model.
    Agreement across gain factors indicates that the learned channel captures the observable-relevant noise components.
    }
    \label{fig:supp-numerical-learning}
\end{figure}

\bibliographystyle{qubit}
\bibliography{ref-codex}

\end{document}
